% Insert into the existing CVAE appendix before the scenario-count sensitivity subsection.
% No additional LaTeX packages required beyond booktabs and graphicx.
\subsection{Comparison with a historical-analogue empirical copula}
\label{app:cvae_copula_96x96}

To evaluate the probabilistic scenario generator against a nonparametric
alternative, we compare the deployed continuous CVAE with a
\emph{historical-analogue empirical-copula benchmark}. Rather than fitting a
parametric copula, the benchmark samples complete paired wind-power and
day-ahead (DA) price trajectories from historical observations, preserving
the empirical temporal and cross-variable rank dependence of each trajectory.
Candidates are restricted to analogue-start months within two calendar months of
the forecast-issue month. Their sampling weights are based on a Gaussian kernel of
historical similarity, with a uniform-weight component to avoid excessive
concentration on a small number of analogues.

The comparison uses \emph{matched recent-information and evaluation windows}:
both predictors are conditioned on the same 96 completed hourly wind-power
and DA-price observations available at the D$-1$ 09:30
\texttt{America/Chicago} forecast issue, and both are scored on the
following 96 consecutive hours, starting immediately after the last completed
observation. For the copula, the similarity features comprise four consecutive
24-hour means for each variable, the most recent 24-hour standard deviation
and the most recent 12-hour mean for each variable; prices are transformed
with $\operatorname{asinh}(\lambda^{\mathrm{DA}}/50)$ for similarity evaluation.
The selected historical DA-price trajectories are additionally recalibrated
using a monotonic empirical-quantile mapping to the last 96 observed DA
prices. This preserves the ranks of the sampled paired trajectories.
The reference library contains complete historical analogue windows ending
no later than June 2025 and is frozen during the 2026 evaluation.

The continuous CVAE retains the checkpoint and paired-residual calibration
used in the MPC case studies. As this model generates 120 future hours from a
96-hour conditioning history, its first 96 forecast hours are used here;
the underlying network is not retrained. Each predictor generates 500 joint
trajectories, reduced by the same backward-deletion algorithm to 50 weighted
representatives, with the reduction distance computed on the common 96-hour
evaluation window. These 96-hour forecasts begin at the first unobserved hour,
not at the delivery-day midnight. Consequently, the exercise is a
\emph{forecast-quality comparison}, distinct from the 84-hour, delivery-aligned
scenario reduction and closed-loop MPC assessment in the main text.

\begin{table}[!htbp]
\centering\small
\caption{Matched 96-hour conditioning/96-hour prediction comparison of the
historical-analogue empirical copula and deployed continuous CVAE on the 2026
unseen holdouts. Each period has four daily forecast issue origins, with
scores averaged over three scenario-sampling seeds (3714, 13714 and 2718).
Both methods generate 500 trajectories reduced to 50 weighted scenarios.
CRPS units are MW for wind and USD/MWh for DA price; the central-interval
coverage target is 90\%, and the trajectory energy score is dimensionless.
Bold indicates the lower CRPS/energy score or coverage closest to 90\% in
each pair. Daily forecast windows overlap and are not independent samples.}
\label{tab:cvae_copula_96x96}
\resizebox{\linewidth}{!}{%
\begin{tabular}{llrrrrr}
\toprule
Period & Predictor & Wind CRPS & DA-price CRPS & Wind cov. (\%) & DA-price cov. (\%) & Energy score \\
\midrule
February 2026 & Copula & \textbf{41.13} & \textbf{4.87} & \textbf{82.3} & 95.4 & \textbf{2.426} \\
              & CVAE & 43.52 & 6.70 & 82.2 & \textbf{93.1} & 2.731 \\
\addlinespace
May 2026 & Copula & 35.75 & 12.09 & \textbf{94.7} & 58.4 & 3.163 \\
         & CVAE & \textbf{35.40} & \textbf{8.73} & 97.1 & \textbf{96.1} & \textbf{2.736} \\
\addlinespace
July 2026 & Copula & 34.65 & 6.88 & \textbf{91.3} & 62.9 & 2.286 \\
          & CVAE & \textbf{32.44} & \textbf{5.15} & 96.4 & \textbf{99.0} & \textbf{1.978} \\
\midrule
Pooled & Copula & 37.17 & 7.95 & \textbf{89.4} & 72.3 & 2.625 \\
       & CVAE & \textbf{37.12} & \textbf{6.86} & 91.9 & \textbf{96.1} & \textbf{2.482} \\
\bottomrule
\end{tabular}}
\end{table}

The copula has lower wind and DA-price CRPS and a lower joint energy score
in February, whereas the CVAE performs better in May and July. The pooled
wind CRPS values are almost identical, while the CVAE achieves lower pooled
DA-price CRPS and joint trajectory energy score across all three sampling
seeds. The copula's 96-observation price-marginal recalibration produces
substantial DA-price interval undercoverage in May and July; the CVAE shows
some overcoverage in these periods. Hence, equal recent-information windows
do not imply equally calibrated models, and the reported comparison does not
establish that the CVAE is uniformly superior or yields greater MPC revenue.
Twelve overlapping forecast issue origins and three Monte Carlo seeds provide
a limited basis for statistical inference. The copula's historical analogue
selection and marginal transformation are design choices for this controlled
benchmark, rather than a claim of an optimally tuned copula model.
